\documentclass[../../main.tex]{subfiles}
\begin{document}

\begin{flushright}
\footnotesize\textit{【自動文字起こし・要確認】}
\end{flushright}

[解] (1) $k=8$の時、 $a_1 = 0, a_2 = 2, a_n = 3 \ (n \ge 3)$ \\
$k=9, \ \ a_1 = 0, a_2 = 3, a_n = 4 \ (n \ge 3)$ \hfill(終)

(2) 題意を帰納的に示す。まず、$n=1$の時、$a_1=0, a_2 = [\dfrac{k}{3}] \ge 0$、$k \ge 0$ から成立する。以下$n=k \in \mathbb{N}$での成立を仮定する。$[x]$は単調増加だから、
\[ a_{k+1} = \left[ \dfrac{a_k + k}{3} \right] \le \left[ \dfrac{a_{k+1} + k}{3} \right] = a_{k+2} \cdots \text{①} \]
次に、
\[ a_{k+1} = \left[ \dfrac{a_k + k}{3} \right] \le \left[ \dfrac{\dfrac{k}{2} - \dfrac{1}{6}}{3} \right] \quad \text{(※修正)} \]
であり、$k \in \text{odd}$の時 $\left[ \dfrac{k}{2} - \dfrac{1}{6} \right] = \left[ \dfrac{k-1}{2} + \dfrac{1}{3} \right] = \dfrac{k-1}{2}$ \\
$k \in \text{even}$の時 $\left[ \dfrac{k}{2} - \dfrac{1}{6} \right] = \dfrac{k}{2} - 1 < \dfrac{k-1}{2}$ \\
から、
\[ a_{k+1} \le \dfrac{k-1}{2} \cdots \text{②} \]
となる。以上①②から、$n=k+1$でも成立。よって示された \hfill(終)

(3) 題意を帰納的に示す。$n=N$の時は成立するので、以下$n=N \in \mathbb{N}$としての成立を仮定する。
\[ a_{N+1} = \left[ \dfrac{a_N + k}{3} \right] = \left[ \dfrac{a_n + k}{3} \right] = a_{n+1} = a_n \quad (\because \text{仮定}) \]
から、$n=N+1$でも成立。よって示された \hfill(終)

(4) (2)から $0 \le a_n \le \dfrac{k-1}{2}$ で、$a_n \in \mathbb{Z}$ から、$a_n = a_{n+1}$ となる$n$がある。この時、
\[ a_{n+1} = \left[ \dfrac{a_n + k}{3} \right] = a_n \]
から、
\[ a_n \le \dfrac{a_n + k}{3} < a_n + 1 \]
\[ \dfrac{k-3}{2} < a_n \le \dfrac{k}{2} \cdots \text{③} \]
である。

$1^\circ \ k \in \text{odd}$ \\
$a_n \in \mathbb{Z}$ と③から、$a_n = \dfrac{k-1}{2}$ である。

$2^\circ \ k \in \text{even}$ \\
$a_n \in \mathbb{Z}$ と③、$a_n \le \dfrac{k-1}{2}$ から、$a_n = \dfrac{k-2}{2}$ である。

以上まとめて、
\[ a_n = \begin{cases} \dfrac{k-1}{2} & (k \in \text{odd}) \\ \dfrac{k}{2} - 1 & (k \in \text{even}) \end{cases} \] \hfill(終)

\end{document}
